\index{BPatch_basicBlockLoop: : : : : : :!}{

An object of this class represents a loop in the code of the application being instrumented. We detect both natural

loops (single-entry loops) and irreducible loops (multi-entry loops). For a natural loop, it has only one entry block

and this entry block dominates all blocks in the loop; thus the entry block is also called the head or the header of

the loop. However, for an irreducible loop, it has multiple entry blocks and none of them dominates all blocks in the

loop; thus there is no head or header for an irreducible loop. The following figure illustrates the difference:}



\begin{tikzpicture}

\begin{scope}[xshift=10.613cm,yshift=-1.173cm]

\path[line width=0.0209in,draw]  (NaN,NaN) --  (Infinity,-Infinity);

\end{scope}

\end{tikzpicture}

\begin{tikzpicture}

\begin{scope}[xshift=9.668cm,yshift=-6.888cm]

\path[line width=0.0102in,draw]  (NaN,NaN) --  (Infinity,NaN) --  (Infinity,-Infinity) --  (NaN,-Infinity) -- cycle;

\node[below,align=justify,text width=NaNcm,transform shape] at (NaN,NaN) {(b) An example of irreducible loop};

\end{scope}

\end{tikzpicture}

\begin{tikzpicture}

\begin{scope}[xshift=12.388cm,yshift=-4.514cm]

\path[line width=0.0209in,draw]  (NaN,NaN) --  (Infinity,-Infinity);

\end{scope}

\end{tikzpicture}

\begin{tikzpicture}

\begin{scope}[xshift=12.829cm,yshift=-2.551cm]

\path[line width=0.0209in,draw]  (NaN,NaN) --  (Infinity,-Infinity);

\end{scope}

\end{tikzpicture}

\begin{tikzpicture}

\begin{scope}[xshift=12.702cm,yshift=-2.937cm]

\path[line width=0.0209in,draw]  (NaN,NaN) --  (Infinity,-Infinity);

\end{scope}

\end{tikzpicture}

\begin{tikzpicture}

\begin{scope}[xshift=11.174cm,yshift=-4.170cm]

\path[line width=0.0209in,draw]  (NaN,NaN) --  (Infinity,-Infinity);

\end{scope}

\end{tikzpicture}

\begin{tikzpicture}

\begin{scope}[xshift=11.026cm,yshift=-2.828cm]

\path[line width=0.0209in,draw]  (NaN,NaN) --  (Infinity,-Infinity);

\end{scope}

\end{tikzpicture}

\begin{tikzpicture}

\begin{scope}[xshift=12.273cm,yshift=-1.141cm]

\path[line width=0.0209in,draw]  (NaN,NaN) --  (Infinity,-Infinity);

\end{scope}

\end{tikzpicture}

\begin{tikzpicture}

\begin{scope}[xshift=2.896cm,yshift=-0.268cm]

\path[line width=0.0102in,draw]  (NaN,-Infinity) --  (Infinity,NaN) --  (Infinity,-Infinity) --  (Infinity,-Infinity) --

cycle;

\node[below,align=center,text width=NaNcm,transform shape] at (Infinity,-Infinity) {Entry};

\end{scope}

\begin{scope}[xshift=1.588cm,yshift=-3.551cm]

\path[line width=0.0102in,draw]  (NaN,-Infinity) -- cycle;

\node[below,align=left,text width=NaNcm,transform shape] at (Infinity,-Infinity) {2};

\end{scope}

\begin{scope}[xshift=3.230cm,yshift=-1.949cm]

\path[line width=0.0102in,draw]  (NaN,-Infinity) -- cycle;

\node[below,align=left,text width=NaNcm,transform shape] at (Infinity,-Infinity) {1};

\end{scope}

\begin{scope}[xshift=4.918cm,yshift=-3.572cm]

\path[line width=0.0102in,draw]  (NaN,-Infinity) -- cycle;

\node[below,align=left,text width=NaNcm,transform shape] at (Infinity,-Infinity) {3};

\end{scope}

\begin{scope}[xshift=2.896cm,yshift=-5.849cm]

\path[line width=0.0102in,draw]  (NaN,-Infinity) --  (Infinity,NaN) --  (Infinity,-Infinity) --  (Infinity,-Infinity) --

cycle;

\node[below,align=center,text width=NaNcm,transform shape] at (Infinity,-Infinity) {Exit};

\end{scope}

\begin{scope}[xshift=3.727cm,yshift=-1.141cm]

\path[line width=0.0209in,draw]  (NaN,NaN) --  (Infinity,-Infinity);

\end{scope}

\begin{scope}[xshift=2.473cm,yshift=-2.838cm]

\path[line width=0.0209in,draw]  (NaN,NaN) --  (Infinity,-Infinity);

\end{scope}

\begin{scope}[xshift=2.618cm,yshift=-4.182cm]

\path[line width=0.0209in,draw]  (NaN,NaN) --  (Infinity,-Infinity);

\end{scope}

\begin{scope}[xshift=4.159cm,yshift=-2.947cm]

\path[line width=0.0209in,draw]  (NaN,NaN) --  (Infinity,-Infinity);

\end{scope}

\begin{scope}[xshift=4.309cm,yshift=-2.582cm]

\path[line width=0.0209in,draw]  (NaN,NaN) --  (Infinity,-Infinity);

\end{scope}

\begin{scope}[xshift=3.842cm,yshift=-4.526cm]

\path[line width=0.0209in,draw]  (NaN,NaN) --  (Infinity,-Infinity);

\end{scope}

\begin{scope}[xshift=1.570cm,yshift=-7.018cm]

\path  (NaN,NaN) --  (Infinity,NaN) --  (Infinity,-Infinity) --  (NaN,-Infinity) -- cycle;

\node[below,align=left,text width=NaNcm,transform shape] at (NaN,NaN) {(a) An example of natural loop};

\end{scope}

\end{tikzpicture}

\begin{tikzpicture}

\begin{scope}[xshift=11.441cm,yshift=-5.835cm]

\path[line width=0.0102in,draw]  (NaN,-Infinity) --  (Infinity,NaN) --  (Infinity,-Infinity) --  (Infinity,-Infinity) --

cycle;

\node[below,align=center,text width=NaNcm,transform shape] at (Infinity,-Infinity) {Exit};

\end{scope}

\end{tikzpicture}

\begin{tikzpicture}

\begin{scope}[xshift=13.458cm,yshift=-3.561cm]

\path[line width=0.0102in,draw]  (NaN,-Infinity) -- cycle;

\node[below,align=justify,text width=NaNcm,transform shape] at (Infinity,-Infinity) {3};

\end{scope}

\end{tikzpicture}

\begin{tikzpicture}

\begin{scope}[xshift=11.774cm,yshift=-1.949cm]

\path[line width=0.0102in,draw]  (NaN,-Infinity) -- cycle;

\node[below,align=justify,text width=NaNcm,transform shape] at (Infinity,-Infinity) {1};

\end{scope}

\end{tikzpicture}

\begin{tikzpicture}

\begin{scope}[xshift=10.142cm,yshift=-3.540cm]

\path[line width=0.0102in,draw]  (NaN,-Infinity) -- cycle;

\node[below,align=justify,text width=NaNcm,transform shape] at (Infinity,-Infinity) {2};

\end{scope}

\end{tikzpicture}

\begin{tikzpicture}

\begin{scope}[xshift=11.441cm,yshift=-0.268cm]

\path[line width=0.0102in,draw]  (NaN,-Infinity) --  (Infinity,NaN) --  (Infinity,-Infinity) --  (Infinity,-Infinity) --

cycle;

\node[below,align=center,text width=NaNcm,transform shape] at (Infinity,-Infinity) {Entry};

\end{scope}

\end{tikzpicture}







\bigskip





\bigskip





\bigskip





\bigskip





\bigskip





\bigskip





\bigskip





\bigskip





\bigskip





\bigskip





\bigskip





\bigskip



{

Figure (a) above shows a natural loop, where block 1 represents the single entry and block 1 is the head of the loop.

Block 1 dominates block 2 and block 3. Figure (b) above shows an irreducible loop, where block 1 and block 2 are the

entries of the loop. Neither block 1 nor block 2 domiantes block 3.}





\bigskip



{

bool containsAddress(unsigned long addr)}



{

Return true if addr is contained within any of the basic blocks that compose this loop, excluding the block of any of

its sub-loops.}



{

bool containsAddressInclusive(unsigned long addr)}



{

Return true if addr is contained within any of the basic blocks that compose this loop, or in the blocks of any of its

sub-loops.}



{

int getBackEdges(std::vector<BPatch_edge *> &edges)}



{

Returns the number of back edges in this loop and adds those edges to the edges vector. An edge is a back edge if it

is from a block in the loop to an entry block of the loop.}



{

int getLoopEntries(std::vector<BPatch_basicBlock *> &entries)}



{

Returns the number of entry blocks of this loop and adds those blocks to the entries vector. An irreducible loop can

have multiple entry blocks.}



{

bool getContainedLoops\index{getContainedLoops: : : : : :

:!}(std::vector<BPatch_basicBlockLoop*>&)}



{

Fill the given vector with a list of the loops nested within this loop.}



{

BPatch_flowGraph *getFlowGraph()}



{

Return a pointer to the control flow graph that contains this loop.}



{

bool getOuterLoops\index{getContainedLoops: : : : : :

:!}(std::vector<BPatch_basicBlockLoop*>&)}



{

Fill the given vector with a list of the outer loops nested within this loop.}



{

bool getLoopBasicBlocks(std::vector<BPatch_basicBlock*>&)\index{getLoopBasicBlocks: : : : : :

:!}}



{

Fill the given vector with a list of all basic blocks that are part of this loop.}



{

bool getLoopBasicBlocksExclusive(}



{

std::vector<BPatch_basicBlock*>&)\index{getLoopBasicBlocks: : : : : : :!}}



{

Fill the given vector with a list of all basic blocks that are part of this loop but not its sub-loops.}



{

bool hasAncestor(BPatch_basicBlockLoop*)\index{getLoopHead: : : : : : :!}}



{

Return true if this loop is nested within the given loop (the given loop is one of its ancestors in the tree of loops).}



{

bool hasBlock(BPatch_basicBlock *b)}



{

Return true if this loop or any of its sub-loops contain b, false otherwise.}



{

bool hasBlockExclusive(BPatch_basicBlock *b)}



{

Return true if this loop, excluding its sub-loops, contains b, false otherwise.}



